% tab:levels
\begin{table}[htbp]
\centering\footnotesize
\caption{GDP per capita in the Baltic states, Finland, the Western core and the USSR, selected years, 1913--2023 (\intlnineteen{})}
\label{tab:levels}
\begin{threeparttable}
\begin{tabular}{lrrrrrr}
\toprule
Year & Latvia & Lithuania & Estonia & Finland & Western & USSR / former \\
& & & & & core & USSR \\
\midrule
1913 & 2,361 & 1,632 & 2,058 & 2,111 & 4,008 & 1,414 \\
1929 & 2,351 & 1,577 & 2,113 & 2,717 & 4,695 & 1,386 \\
1938 & 3,099 & 2,057 & 2,685 & 3,589 & 5,164 & 2,150 \\
\midrule
1950 & 3,722 & 2,298 & 3,667 & 4,253 & 5,569 & 2,586 \\
1968 & 7,544 & 5,399 & 7,348 & 8,093 & 10,542 & 4,735 \\
1973 & 8,523 & 6,464 & 8,080 & 11,085 & 12,771 & 5,515 \\
1989 & 8,772 & 7,858 & 8,958 & 16,946 & 17,329 & 6,473 \\
\midrule
1995 & 5,018 & 5,203 & 6,277 & 16,162 & 18,769 & 3,933 \\
2008 & 13,015 & 13,099 & 14,604 & 25,175 & 25,157 & 9,266 \\
2019 & 16,222 & 18,184 & 17,883 & 25,016 & 27,652 & 11,060 \\
2023 & 17,603 & 19,162 & 17,837 & 25,142 & 27,595 & 11,482 \\
\bottomrule
\end{tabular}
\begin{tablenotes}[para,flushleft]
\footnotesize
\textit{Notes:} USSR within 1990 borders, former Soviet Union after 1991; from 1946 the Soviet level is Maddison's series rescaled to Bolotin's 1990 benchmark, to which the Baltic series are anchored. Western core as in Table~\ref{tab:growth}; the annual Latvian series is in Table~\ref{tab:series}. \textit{Sources:} as Table~\ref{tab:growth}.
\end{tablenotes}
\end{threeparttable}
\end{table}
